# Profil E2 — Prise de notes brute (Associate Partner, 25 ans)

**Code :** E2  
**Fonction :** Associate Partner  
**Ancienneté :** 25 ans  
**Contexte :** Institutions européennes (BEI). Développement commercial, pitchs. Sponsor practice Data & AI Luxembourg.

---

## T0 — Cadrage

**Q : Parcours, comment arrivé ici ?**
- 25 ans. Associé.
- BEI. Développement commercial. Pitchs. Cherche nouveaux clients.
- Supervise 4 juniors.
- Vu évolution conseil. Papier à IA générative.

---

## T1 — Usage réel

**Q : Dernière fois utilisé outil génératif ?**
- Ce matin. DeepSearch. Recherche Sharepoints.
- Projet similaire. Superprompt précis. Tourné 25 min.
- Synthèse parfaite. Des heures manuellement.
- Accélérateur de compétence, pas substitut.

**Q : Sur quelles tâches au quotidien ?**
- Tout. Pitchs structures/angles. Livrables premiers jets. Veille concurrentielle.
- DeepSearch efficace.
- Jamais données confidentielles dans outils non-officiels. Vérifie tout.
- Outil, pas fin.

**Q : Comment appris ?**
- Autodidacte. Sponsor Data & AI. Suis presse. Teste nouveautés.
- Partage prompts qui marchent.
- Passionne. Potentiel énorme.

**Q : Temps gagné mis sur quoi ?**
- Relation client. Recherche/rédaction/structuration -> réinvesti écoute, compréhension besoins.
- Cœur métier = comprendre, pas produire.
- IA permet mieux comprendre. Plus temps pour creuser, challenger client. Vraie valeur ajoutée.

**Q : Situation où plus d'effort que de le faire soi-même ?**
- Rare. Sais quand utiliser. Sujets spécifiques/techniques/données sensibles = direct.
- 80% du travail gain net. Savoir où pertinent, où pas.

**Q : Vous en parlez entre vous ?**
- Beaucoup. Certains réticents. Peur contrôle, déshumanisation.
- Moi pense inverse. IA rend plus pertinent.
- Débat pas fini. Techno rupture prend temps. Mon rôle = accompagner transition.

---

## T2 — Apprentissage

**Q : Comment on apprend le métier aujourd'hui ?**
- Différemment. Avant : heures chercher, structurer, rédiger. Junior fait fraction temps.
- Mais doit apprendre autre chose : vérifier, critiquer, s'assurer pertinence.
- Déplacement compétences, pas disparition.
- Bon consultant = sait poser bonnes questions, évaluer réponses. Compétence nouvelle à enseigner.

**Q : Par quelles tâches on apprend vraiment ?**
- Contact client. IA produit contenu, pas sentir hésitation, non-dit, culture.
- Dis juniors : "IA pour gagner temps, mais utilise temps terrain, échanger, observer." Comme ça bon consultant.

**Q : Tâches ingrates ?**
- Formatives, mais pas même proportion. Juniors doivent faire 1x manuel pour comprendre difficulté.
- 1x compris, encourage IA pour aller vite. Important pas sauter compréhension.
- Déléguer tout IA sans jamais fait = sauras pas évaluer.

**Q : Comment vérifiez ce que produit junior qui utilise IA ?**
- Pas différemment. Qualité, cohérence, pertinence.
- IA ou pas change pas fond. Demande : "qu'est-ce apporté ?"
- Copie IA = coquille vide. Jugement/compréhension/nuance = ce que veux voir. Outil, pas remplaçant.

---

## T3 — Client et valeur

**Q : Confiance client repose sur quoi ?**
- Être en avance. Anticiper besoins, proposer solutions pas imaginées.
- IA aide en avance. Plus temps réfléchir, angles pas vus.
- Client paie vision, pas temps production. IA bien utilisée amplifie vision.

**Q : Sujet IA déjà abordé avec client ?**
- Très souvent. J'aborde volontairement.
- Dis : "IA pour être + pertinents, + rapides, + créatifs. Partenaire, pas substitut."
- Clients apprécient transparence. Veulent savoir à pointe. Argument commercial. Pas cacher.

**Q : Diriez-vous au client qu'une partie est IA ?**
- Oui. "IA pour produire, retravaillé adapté besoins." Client comprend. Utilise peut-être IA. Qualité produit final compte, pas méthode.

**Q : Livrable moins crédible ?**
- Manque personnalisation. IA générique. Si client sent copié-collé sans adapter, doute.
- Problème jugement, pas technologie.
- Bon consultant sait quand IA utile, quand tout reprendre. Crédibilité = valeur ajoutée où IA pas peut aller.

**Q : Dans 10 ans, client paiera pour quoi ?**
- Intuition et relation. IA fera beaucoup, pas remplacer relation humaine.
- Consultant = "traducteur" entre IA et client. Interprète IA, rend pertinent. Rôle nouveau, passionnant.

---

## T4 — Gouvernance

**Q : Règles sur usage ?**
- Politique. Trop timide. Autorise Copilot, interdit autres. Pas dit comment vérifier, assurer qualité.
- Sponsor Data & AI = aider construire règles. Pas freiner, sécuriser.
- Veux gens utilisent IA en connaissance de cause. Règle claire = protection, pas entrave.

**Q : Avant livrable client ?**
- Revue. Pas diff entre IA/non-IA. Vérifie fond, cohérence, pertinence.
- Consultant justifie ce qu'il met. IA pas excuse. Aide. Demande "pourquoi" pas "qu'est-ce".

**Q : Déjà rencontré contenu faux IA ?**
- Oui. Pas plus qu'erreurs humaines. IA peut halluciner. Consultant peut se tromper.
- Erreur IA + facile détecter si sais chercher. Forme équipes vérification critique. Compétence, pas défiance.

**Q : Règle idéale ?**
- "IA outil amplification. Utilisez pour aller + loin, pas -."
- Formez vérification, critique, personnalisation. Co-pilote, pas pilote auto.
- N'ayez pas peur erreurs. Apprenez les repérer. Comme ça devient meilleur.

---

## T5 — Clôture

**Q : À la place de la direction, première décision ?**
- Initiative partage bonnes pratiques bureau. Pas formation descendante, espace échanges.
- Ressources temps/moyens pour expérimenter. IA s'apprend en pratiquant, pas lisant charte. Direction = créer conditions expérimentation.

**Q : Aspect important non abordé ?**
- Image cabinet. Identifié cabinet qui maîtrise IA = avantage concurrentiel.
- Clients regardent. Veulent voir aligné sur leur transformation. Réputation, pas juste efficacité. Montrer à pointe via projets concrets, cas d'usage, succès.

**Q : Point de vue utile à recueillir ?**
- Responsable juridique. Voir comment voient risques : confidentialité, PI, responsabilité.
- Avocats prudents, pessimistes. Leur regard important équilibrer enthousiasme. Dialogue à avoir.